% Title page, revision history, how to read, contents, list of figures.
\begin{titlepage}
  \centering
  \vspace*{3cm}
  {\Huge\bfseries Software Design Document\par}
  \vspace{0.6cm}
  {\LARGE for\par}
  \vspace{0.6cm}
  {\Huge\bfseries \SRSproduct\par}
  \vspace{0.4cm}
  {\Large Design of Release 1 against the Software Requirements Specification, version \SRSversion\par}
  \vspace{1.5cm}
  {\Large Version \SDDversion\par}
  \vspace{0.4cm}
  {\large \SDDdate\par}
  \vfill
  {\large Prepared by\par}
  {\Large Jai Nana \quad Brian Lujan\par Fotih Turopov \quad Yerassyl Shnazbayev\par}
  \vspace{0.8cm}
  {\large CSDS 393 Software Engineering\par Case Western Reserve University\par}
  \vspace{1cm}
\end{titlepage}

\section*{Revision History}
\begin{tabularx}{\textwidth}{@{}L{3.2cm}L{2.2cm}YL{2.6cm}L{1.4cm}@{}}
\toprule
\textbf{Name} & \textbf{Date} & \textbf{Reason for changes} & \textbf{Reviewed by} & \textbf{Version} \\
\midrule
J.\ Nana & \SDDdate & First complete draft: architecture, module specifications, runtime scenarios, data, API, crosscutting concerns, evolution, traceability. & Team inspection, Appendix~\ref{app:inspection} & 1.0 \\
\bottomrule
\end{tabularx}

\noindent\textbf{How to read this document.} Section~\ref{sec:intro} states the baseline and the principles; Section~\ref{sec:arch} gives the architecture in four figures; Section~\ref{sec:decomp} allocates every responsibility to a module; Section~\ref{sec:modules} specifies each module with one template; Section~\ref{sec:runtime} shows the transactions, the state machines, and seven scenarios as sequence diagrams; Sections~\ref{sec:data} to \ref{sec:cross} give the schema, the API, and the crosscutting concerns; Section~\ref{sec:evo} says how the design evolves. Appendix~\ref{app:trace} traces every requirement of the SRS to a module and a section. Capitalized words such as Request, Hold, and Loan have the meaning given in the glossary, Appendix~\ref{app:glossary}; code identifiers are set in \code{monospace}; module names in \module{bold sans-serif}.

\tableofcontents
\listoffigures
\clearpage
